\section{Background on Euclidean projections} \label{sec:euclidean_projections}

At the core of many problems in optimisation, statistics, and machine learning lies the task of identifying a point that satisfies multiple constraints simultaneously~\cite[Ch. 8]{BOYDCONVEX}. Formally, we consider a Euclidean space $\mathcal{X} \subseteq \Rset^p$ equipped with the Euclidean norm $\twonorm{x}\eqdef (x_1^2+\dots+x_p^2)^{1/2}$, and a finite family of $n$ closed convex subsets $\{\mathcal{H}_i\}_{i=0}^{n-1}\subseteq\mathcal{X}$ with a nonempty intersection $\mathcal{H} := \bigcap_{i=0}^{n-1} \mathcal{H}_i$. Given a point $x^\circ \in \mathcal{X}$, its Euclidean projection $x^\star \reqdef \mathcal{P}_{\mathcal{H}}(x^\circ)$ onto set $\mathcal{H}$ can be found by solving the \textit{constrained quadratic program} (CQP)
\begin{equation} \label{eq:projection}
\begin{array}{ll}
    {\displaystyle \minimise_{x\in\mathcal{X}}}  & \frac{1}{2}\twonorm{x-x^\circ}^2 \\
    \text{subject to} & x \in \bigcap_{i=0}^{n-1}\mathcal{H}_i,
\end{array}
\end{equation}
where $x$ is the decision variable, and $x^\circ$ and sets $\{\mathcal{H}_i\}_{i=0}^{n-1}$ are the problem data. In this paper, we assume $\mathcal{H}$ to be a polyhedral set, meaning $\mathcal{H}=\set{x \in \mathcal{X}}{A x \leq b}$, with $A \in \Rset^{n\times p}$ and $b \in \Rset^n$, and each $\mathcal{H}_i$ a half-space, written as
\begin{align}\label{eq:sets_i}
\mathcal{H}_i\coloneqq\set{x\in\mathcal{X}}{a_i^Tx \leq b_i}, \quad i=0,\dots,n-1,
\end{align}
where vectors $a_i$ are unit normals to each bounding hyperplane $\partial\mathcal{H}_i\coloneqq\set{x\in\mathcal{X}}{a_i^Tx = b_i}$.

The formulation in~\eqref{eq:projection} is widely adopted in optimisation applications, ranging from restricted least squares and isotonic regression in statistics, to signal and image processing (where convex sets encode prior knowledge, like non-negativity or sparsity), control engineering (where system constraints can be modelled using convex sets), and machine learning (where solutions to learning problems must often adhere to structural constraints~\cite{BAUSCHKESWISS}). While in some cases the projection onto $\mathcal{H}$ is easy to compute, for arbitrary polyhedral sets the projection operator $\mathcal{P}_{\mathcal{H}}$ onto the intersection set is intractable and no closed-form solution to~\eqref{eq:projection} is known. In these cases, the solution to~\eqref{eq:projection} can be obtained using standard optimisation software packages. To solve~\eqref{eq:projection}, most solvers require introducing an additional variable $z\eqdef Ax$ (e.g., \cite{osqp} or~\cite{ocpb:16}) that is projected onto the set $\set{z\in\Rset^n}{z \leq b}$, for which an explicit solution exists~\cite{BOYDCONVEX}. However, for large $n$, this variable augmentation can reduce the computational efficiency, particularly in high-dimensional or time-sensitive applications when the projection is part of an overarching algorithm~\cite{OPTIMNESTEROV,KEMPF20206542}.

As an alternative that can be faster in practice~\cite{KEMPF20206542}, \textit{Dykstra's projection algorithm}~\cite{DYKSTRA} is an iterative method designed to solve~\eqref{eq:projection}. It is an extension of the \textit{method of alternating projections} (MAP), which, given an initial iterate, finds a point in the intersection of the half-spaces. Dykstra's algorithm modifies MAP by introducing auxiliary variables to achieve strong convergence to the true projection. Qualitatively, Dykstra's algorithm stores the negative direction vector of projection into auxiliary variables; these are added to the primary variables before subsequent projections and, as a result, the primary iterates of Dykstra's method asymptotically converge to the optimal solution of~\eqref{eq:projection}. Although, as in standard optimisation packages, Dykstra's method introduces additional variables, these appear only in vector additions rather than in matrix-vector multiplications. This reduces the computational load and can make Dykstra's method a highly valuable algorithm to deploy in real-time, high-throughput applications for finding fast (approximate) solutions to CQPs. However, despite its strong theoretical guarantees, Dykstra's algorithm is known to suffer from a \emph{stalling} problem~\cite{DYKSTRASTALLING}. In such cases, the iterates of Dykstra's method remain unchanged for a number of iterations. The duration of the stalling period is not known \emph{a priori}, and, given the choice of starting point, can be made arbitrarily long. In general, the algorithm cannot be run for a fixed number of iterations with a guarantee that the output will be closer to the projection than the initial point, which can hinder deployment of Dykstra's method in practical applications.

In this paper, we address the stalling problem of Dykstra's method for polyhedral sets. Exploiting simplifications that arise in the polyhedral case, we derive the exact duration of the stalling period once stalling is detected (Theorem~\ref{thm:nstall}). This quantity enables a modified version of Dykstra's algorithm (Algorithm~\ref{alg:fastforward}) that fast-forwards through the stalling period, thereby improving convergence properties. Crucially, this modification preserves all convergence guarantees of the original algorithm (Corollary~\ref{cor:convergence}). This development represents a step towards practical deployment of Dykstra's algorithm in real-time settings, where its low-overhead projection could offer significant performance advantages. The structure of the paper is as follows. Section~\ref{sec:background} reviews MAP, Dykstra's algorithm, and narrows the analysis to polyhedral sets, introducing the stalling phenomenon. Section~\ref{sec:main_result} proposes the fast-forward modification, and Section~\ref{sec:example} presents supporting numerical experiments, which demonstrate the effectiveness of the solution. Finally, Section~\ref{sec:conclusion} concludes with a discussion of potential applications, extensions, and refinements. 


\emph{Notation}.
We denote the transpose of a vector $a$ as $a^T$, and its $\ell_{2}$-norm as $\twonorm{a}$. The interior of a set $\mathcal{X}$ is denoted as $\text{int}\,\mathcal{X}$, and its boundary is denoted as $\partial\mathcal{X}$. Function composition is written as $g \circ h$, and the modulo and ceiling operators are denoted by $[\cdot]$ and $\lceil \cdot \rceil$, respectively. The absolute iteration index is denoted as $m$, the half-space index (i.e. which half-space we are currently considering) is denoted as $i$, while the half-space count (i.e. how many half-spaces there are in total) is denoted as $n$, and is fixed. We name an update with an increase by $1$ in $m$ ($m \leftarrow m+1$) an \emph{iteration}, and refer to an advancement of $n$ iterations ($m \leftarrow m+n)$ as a \emph{cycle}.